\subsection{Complete compiler bounds from commuting native factors}

The homogeneous kinetic and analyzer source decomposes into $6s$ disjoint
four-mode fermionic factors on $s$ sites. Fourier transformation pairs
opposite momenta, and invariance of the native spin epsilon tensor gives
two helicity pairings per momentum and character. Every complementary
mode has identical free evolution in the actual and target laws. On each
factor the rank-one pairing has norm one. Centering its kinetic number
operator gives $\|H_r-2\kappa I\|\le 2|\kappa|+1\le4/a+1$.
The identical scalar phases cancel in the unitary comparison without
changing the recorded energy or reference.

Even operators on disjoint fermionic factors commute. Applying the same
exact fourth-order word to each factor and telescoping the tensor product
gives the complete unhalved trace bound
\[
 \epsilon\le\min\left(2,\frac{6sR^5|t|^5}{10n^4}\right),
 \qquad R=\frac73(4/a+1),\qquad n\ge R|t|.
\]
It includes all ancillary outputs and arbitrary quantum references.
A coupling between factors invalidates the factorization and is retained
as a failing control. The transformation proves an identity of the
existing pulses; it adds no Fourier apparatus operation.

For a nonnegative physical excitation observable with norm at most
$E_{\max}$, a target mean at most $\mu$ and an isometry error $\delta$,
purification gives
\[
 |\langle A\rangle_{\rm actual}-\langle A\rangle_{\rm target}|
 \le2\sqrt{E_{\max}\mu}\,\delta+E_{\max}\delta^2.
\]
The compact near-rest packet has $E_{\max}=282s$ and
$\mu\le m_{\rm upper}+d$. Retaining its energy error budget $d/10^6$
and complete-record trace tolerance $10^{-12}$ requires at most
$6.564\,10^{27}$ cycles under these bounds, a reduction by approximately
$1.966\,10^{34}$ from the unseparated global-norm estimate.
This is an analytic resource bound on the same huge carrier. Coherent
control services, preparation costs and the physical-energy-marginal
domain are explicit; an efficient local or autonomous implementation is
not supplied.
